% Chapter 6: Equivariant Quantum Cohomology of Flag Manifolds
% Based on Bumsig Kim (1996) \cite{Kim1996}
% "On Equivariant Quantum Cohomology"
% arXiv: alg-geom/9608001

\chapter{Equivariant Quantum Cohomology of Flag Manifolds}
\label{chap:EquivariantQuantum}

\section{Historical Context and Motivation}

The classical enumerative geometry of complex projective spaces asks: given $2k+1$ generic points in $\mathbb{CP}^k$, how many degree-$d$ rational curves pass through them? For $k=1$ (curves in $\mathbb{CP}^2$), the answer is $d^2$. Such enumerative problems motivated the development of Gromov-Witten invariants and quantum cohomology.

In 1995, Givental and Kim \cite{GiventalKim1995} established a deep connection between the quantum cohomology of complete flag manifolds and the Toda lattice, showing that the quantum cohomology ring is isomorphic to the coordinate ring of the Toda lattice phase space. Kim's 1996 paper \cite{Kim1996} generalizes this to the equivariant setting, establishing a systematic framework for computing $U_n$-equivariant quantum cohomology of partial flag manifolds.

The key insight is that the equivariant theory specializes to both the ordinary quantum cohomology (by setting $c_1 = \cdots = c_n = 0$) and the usual equivariant cohomology (by setting $q = 0$).

\section{Flag Manifolds and Tautological Bundles}

Let $F_{s_0, s_1, \ldots, s_l}$ denote the manifold of all flags
\begin{equation}
0 \subset \mathbb{C}^{s_0} \subset \mathbb{C}^{s_1} \subset \cdots \subset \mathbb{C}^{s_l} = \mathbb{C}^n
\end{equation}
of complex linear subspaces in $\mathbb{C}^n$ of dimensions $0 < s_0 < s_1 < \cdots < s_l = n$.

\begin{Definition}[\textbf{Tautological bundles}]
\label{def:TautologicalBundles}
For each $i = 0, 1, \ldots, l$, let $k_i$ denote the dimension of the $i$-th tautological bundle:
\begin{equation}
k_0 = s_0, \quad k_1 = s_1 - s_0, \quad \ldots, \quad k_l = s_l - s_{l-1}.
\end{equation}
The $i$-th tautological bundle $\mathcal{V}_i$ has fiber $\mathbb{C}^{s_i}$ over the flag, and the successive quotients $\mathcal{V}_i / \mathcal{V}_{i-1}$ have fibers $\mathbb{C}^{k_i}$.
\end{Definition}

The Chern polynomial of an $m$-dimensional complex vector bundle with Chern classes $c_1, \ldots, c_m$ is:
\begin{equation}
x^m + c_1 x^{m-1} + \cdots + c_m.
\end{equation}

Denote by $P_0, P_1, \ldots, P_l$ the Chern polynomials of the tautological bundles $\mathcal{V}_0, \mathcal{V}_1/\mathcal{V}_0, \ldots, \mathcal{V}_l/\mathcal{V}_{l-1}$ respectively. The cohomology algebra $H^*(F_{s_0, \ldots, s_l}, \mathbb{Q})$ is multiplicatively generated by the Chern classes $c_j^{(i)}$ of these bundles, with a complete set of relations given by:
\begin{equation}
P_0(x) \cdot P_1(x) \cdot \ldots \cdot P_l(x) = x^n.
\tag{1}
\end{equation}

\section{Quantum Multiplication}

The quantum cohomology algebra of a compact K\"ahler manifold $X$ is, by definition, the cohomology space equipped with a new multiplication $\star$. For cohomology classes $a, b, c$ represented by cycles Poincar\'e-dual to them, the structural constant $\langle a \star b, c \rangle$ of the quantum multiplication counts holomorphic maps $\mathbb{CP}^1 \to X$ with marked points mapped to cycles $a, b, c$, weighted by $q^d$ where $d = (d_1, \ldots, d_l) \in H_2(X, \mathbb{Z})$ is the degree and $q^d = q_1^{d_1} \cdots q_l^{d_l}$.

\begin{Remark}
When $q = 0$, the quantum multiplication reduces to the ordinary cup-product:
\[
\langle a \star b, c \rangle|_{q=0} = \langle a \cdot b, c \rangle.
\]
\end{Remark}

The parameters $q_i$ correspond to the generators of $H_2(X, \mathbb{Z})$ represented by rational curves in the flag manifold. Specifically, if $p_i = c_1^{(i)} + \cdots + c_1^{(l)}$ denotes the first Chern class of the quotient bundle, then $q_i^{d_i}$ represents holomorphic curves $C$ with $\int_C p_i = d_i$.

\section{The Connected Fraction Formula}

The key innovation in \cite{Kim1996} is the connected fraction formula for the quantum deformation of relation \eqref{1}. Consider the continued fraction:
\begin{equation}
P_0 + \frac{(-1)^{k_0+1} q_1}{P_1 + \frac{(-1)^{k_1+1} q_2}{P_2 + \frac{\cdots}{P_l}}}.
\tag{2}
\end{equation}

This fraction can be written unambiguously as the ratio $P(x)/Q(x)$ of polynomials of degree $n$ and $n - s_0$, respectively. The numerator takes the form:
\begin{equation}
P(x) = x^n + \Sigma_1 x^{n-1} + \cdots + \Sigma_n,
\tag{3}
\end{equation}
where $\Sigma_j$ are polynomial expressions in the Chern classes $c_j^{(i)}$ and the quantum parameters $q_1, \ldots, q_l$.

\begin{Theorem}[\cite{Kim1996} Theorem 1]
\label{th:QuantumCohomologyFlag}
The quantum cohomology algebra of the flag manifold $F_{s_0, \ldots, s_l}$ is multiplicatively generated by the $n$ Chern classes $c_j^{(i)}$ and the $l$ parameters $q_1, \ldots, q_l$, satisfying the relations
\[
\Sigma_1 = 0, \quad \ldots, \quad \Sigma_n = 0.
\]

The Poincar\'e intersection pairing $(a, b)$ between two cohomology classes represented by polynomials $a(c, q), b(c, q)$ is given by the $n$-dimensional residue
\begin{equation}
\langle a, b \rangle(q) = \left(\frac{1}{2\pi\sqrt{-1}}\right)^n \oint_{| \Sigma_i | = \epsilon_i} a(c, q) b(c, q) \frac{d c_1^{(1)} \wedge \ldots \wedge d c_{k_l}^{(l)}}{\Sigma_1(c, q) \cdot \ldots \cdot \Sigma_n(c, q)}.
\tag{4}
\end{equation}
\end{Theorem}

\begin{Corollary}[\cite{Kim1996} Corollary to Theorem 1]
\label{cor:EnumerativeCounting}
For $N$ given generic cycles $a_1, \ldots, a_N$ of real codimension 2 in the flag manifold, the number of degree $d$ holomorphic maps
\[
(\mathbb{CP}^1, p_1, \ldots, p_N) \to (F_{s_0, \ldots, s_l}, a_1, \ldots, a_N)
\]
is equal to the residue \eqref{4} with $a = a_1(c) a_2(c) \cdots a_N(c)$ and $b = 1$.
\end{Corollary}

\section{The Equivariant Generalization}

Let $U_n$ act on $\mathbb{C}^n$ by left multiplication. This induces an action on the flag manifold $F_{s_0, \ldots, s_l}$. The $U_n$-equivariant cohomology of $X$ is defined as
\[
H_{U_n}^*(X) = H^*(X_{U_n}),
\]
where $X_{U_n} = X \times_{U_n} EU_n$ is the associated bundle over $BU_n$.

Denote by $c_1, \ldots, c_n$ the universal Chern classes of the principal $U_n$-bundle.

\begin{Theorem}[\cite{Kim1996} Theorem 2]
\label{th:EquivariantQuantumCohomology}
The $U_n$-equivariant quantum cohomology $\mathbb{Q}[\nu_1, \ldots, \nu_n]$-algebra of the flag manifold $F_{s_0, \ldots, s_l}$ is isomorphic to
\[
\mathbb{Q}[\Sigma_1^{(1)}, \ldots, \Sigma_{s_l}^{(l)}, \Sigma_{11}, \ldots, \Sigma_{1n}, \Sigma_1, \ldots, \Sigma_k] / (\Sigma_{11} - \Sigma_1, \ldots, \Sigma_{1k} - \Sigma_k),
\]
where:
\begin{itemize}
\item $\Sigma_j^{(i)}$ are the coefficients from the Chern polynomials $P_i(x) = x^{k_i} + \Sigma_1^{(i)} x^{k_i-1} + \cdots + \Sigma_{k_i}^{(i)}$ of the tautological bundles,
\item $\Sigma_{1i}$ denotes the power sum $c_1^i + c_2^i + \cdots + c_n^i$ of the universal Chern classes, related to the elementary symmetric polynomials via Newton's identities,
\item the relations $\Sigma_{1i} - \Sigma_i = 0$ connect the equivariant generators to the ordinary ones.
\end{itemize}

The equivariant Poincar\'e pairing is given by the residue
\begin{equation}
(a, b)(q, c_1, \ldots, c_n) = \left(\frac{1}{2\pi\sqrt{-1}}\right)^n \oint_{| \Sigma_i - c_i | = \epsilon_i} ab \frac{\bigwedge_{i=0}^l \bigwedge_{j=1}^{k_i} d c_j^{(i)}}{\prod_{m=1}^n (\Sigma_m(c, q) - c_m)}.
\tag{5}
\end{equation}
\end{Theorem}

\begin{Remark}
The original paper contains transcription errors in the statement (the notation $\mathbb{Q}[\nu, \ldots, \nu]$ is clearly corrupted). The interpretation above follows the mathematical content: the ring is generated by the Chern class generators of the tautological bundles together with the universal Chern classes of $U_n$, subject to relations that specialize to the non-equivariant case when $c_1 = \cdots = c_n = 0$.
\end{Remark}

The key property is that by setting $c_1 = \cdots = c_n = 0$, we recover the non-equivariant quantum cohomology of \cref{th:QuantumCohomologyFlag}. By setting $q = 0$, we obtain the usual equivariant cohomology.

\section{The Three Fundamental Properties}

Kim establishes three fundamental properties of equivariant quantum cohomology that allow computation and specialization between different settings.

\begin{Lemma}[\cite{Kim1996} Lemma, Properties of Equivariant Quantum Cohomology]
\label{def:ThreeProperties}
Let $G$ be a connected compact Lie group acting on a generalized flag variety $X$. Then there is a $\mathbb{Z}$-graded equivariant quantum cohomology algebra $QH_G^*(X, \mathbb{Q})$ that is $H_G^*(X, \mathbb{Q}) \otimes_{\mathbb{Q}} \mathbb{Q}[q]$ as a free $H^*(BG, \mathbb{Q}) \otimes_{\mathbb{Q}} \mathbb{Q}[q]$-module. The algebra has the following properties:

\subsection{Product Rule}
Let $G'$ and $G''$ be connected compact Lie groups with actions on $X'$ and $X''$ respectively. Then
\begin{equation}
QH_{G' \times G''}^*(X' \times X'', \mathbb{Q}) \cong QH_{G'}^*(X', \mathbb{Q}) \otimes_{\mathbb{Q}} QH_{G''}^*(X'', \mathbb{Q}).
\tag{6}
\end{equation}

\subsection{Restriction Rule}
Let $G'$ be a connected Lie subgroup of a connected compact Lie group $G$, with $X$ a $G$-space. Then, as $H^*(BG', \mathbb{Q})$-algebras,
\begin{equation}
QH_{G'}^*(X, \mathbb{Q}) \cong QH_G^*(X, \mathbb{Q}) \otimes_{H^*(BG)} H^*(BG').
\tag{7}
\end{equation}

\subsection{Induction Rule}
Let $G'$ be a connected Lie subgroup of $G$, and let $G'$ act on $Y$. Define $X := G \times_{G'} Y$, which has the induced $G$-action. If $X$ becomes a generalized flag manifold with holomorphic quotient maps $X \to Y$ and $X \to G/G'$, then, as $H^*(BG, \mathbb{Q})$-algebras,
\begin{equation}
QH_G^*(X, \mathbb{Q}) \otimes_{\mathbb{Q}[q, q']} \mathbb{Q}[q] \cong QH_{G'}^*(Y, \mathbb{Q}),
\tag{8}
\end{equation}
where $q$ (resp. $q'$) is the formal multi-variable for a suitable basis of $H_2(Y, \mathbb{Z})$ (resp. $H_2(G/G', \mathbb{Z})$).
\end{Lemma}

These properties are proved in Section 4 of \cite{Kim1996} using the geometry of stable maps and the fibration $X_{G'} \to BG'$.

\begin{Remark}[Explanation of the Suitable Basis]
The suitable basis of $H_2(X, \mathbb{Z})$ consists of elements represented by rational curves in $X$. Let $X = G/P$ where $G$ is a complex Lie group and $P$ is a parabolic subgroup containing a Borel subgroup $B$. Let $P'$ be a parabolic subgroup containing $P$ with exactly one more root than $P$. Then the fibers of $G/P \to G/P'$ are rational curves, each representing an element of $H_2(X, \mathbb{Z})$. Varying the choice of $P'$ among all such parabolic subgroups provides the basis.
\end{Remark}

\begin{Remark}[Degenerating Leray Spectral Sequence]
In the induction setting $X = G \times_{G'} Y$, the Leray spectral sequence of the fibration $X \to Y$ degenerates, giving:
\[
H_2(X, \mathbb{Z}) \cong H_2(Y, \mathbb{Z}) \oplus H_2(G/G', \mathbb{Z}).
\]
Under this identification, the suitable basis of $H_2(X)$ decomposes into the suitable basis of $H_2(Y)$ and a basis of $H_2(G/G')$. This decomposition explains the appearance of the two formal variables $q$ and $q'$ in the induction formula \eqref{8}.
\end{Remark}

\section{From Theorem 2 to Theorem 1: Specialization}

A key achievement of \cite{Kim1996} is deducing \cref{th:QuantumCohomologyFlag} from \cref{th:EquivariantQuantumCohomology} via two fundamental specializations:

\subsection{Specialization (a): Non-Equivariant Quantum Cohomology}
The restriction rule \eqref{7} shows that when we mod out by the $G$-characteristic classes (i.e., set $c_1 = \cdots = c_n = 0$), the equivariant quantum cohomology becomes the ordinary non-equivariant quantum cohomology:
\[
QH_{U_n}^*(F_{s_0,\ldots,s_l}) \otimes_{H^*(BU_n)} \mathbb{Q} \cong QH^*(F_{s_0,\ldots,s_l}).
\]
This is because the restriction map $H^*(BU_n) \to \mathbb{Q}$ sends all Chern classes to zero.

\subsection{Specialization (b): Usual Equivariant Cohomology}
When we set $q = 0$ (i.e., consider only degree-zero maps, which are constant maps), the quantum multiplication reduces to the classical multiplication. The definition of equivariant quantum cohomology algebras given in Subsection 3.3 shows that:
\[
QH_G^*(X) \|_{q=0} = H_G^*(X),
\]
the ordinary equivariant cohomology.

Combining these two specializations with the induction and product rules, Kim proves Theorem 1 as the specialization of Theorem 2 at $c_1 = \cdots = c_n = 0$. The connected fraction formula \eqref{2} arises from applying these rules to reduce the computation to Grassmannians, where the quantum cohomology was already known (Siebert-Tian, Witten).

\subsection{The Strategy for Flag Manifolds}
Kim's strategy for computing the quantum cohomology of general flag manifolds is:
\begin{enumerate}
\item Use the product rule to reduce products of flag manifolds to individual factors.
\item Use the induction rule to build up flag manifolds from simpler ones (Grassmannians).
\item Use the restriction rule to relate equivariant and non-equivariant quantum cohomology.
\item For Grassmannians $F_{k,n}$, use the known results of Witten and Siebert-Tian.
\end{enumerate}
This hierarchy of computations, combined with the explicit formulas for the tautological bundles, yields the connected fraction formula.

\section{Gromov-Witten Classes}

The foundation of quantum cohomology is the theory of Gromov-Witten classes. Let $\overline{\mathcal{M}}_n(X, d)$ denote the moduli space of stable maps of degree $d \in H_2(X, \mathbb{Z})$ from a genus zero curve with $n$ ordered marked points to $X$.

\begin{Definition}[\cite{Kim1996} Definition in Section 3.1]
\label{def:GromovWittenClass}
For a generalized flag variety $X$, the Gromov-Witten class
\[
I_{n,d}^X: H^*(X)^{\otimes n} \to H^*(\overline{\mathcal{M}}_n(X, d))
\]
is defined by
\begin{equation}
I_{n,d}^X(a_1 \otimes \cdots \otimes a_n) := (\pi^X)_* (ev^X)^*(a_1 \otimes \cdots \otimes a_n),
\tag{9}
\end{equation}
where $\pi^X: \overline{\mathcal{M}}_n(X, d) \to \overline{\mathcal{M}}_n$ and $ev^X: \overline{\mathcal{M}}_n(X, d) \to X^n$ are the projection and evaluation maps respectively.
\end{Definition}

The moduli space $\overline{\mathcal{M}}_n(X, d)$ has the ``right'' dimension (as noted in \cite{Kim1996}):
\[
\dim \overline{\mathcal{M}}_n(X, d) = \int_d c_1(T_X) + \dim X + n - 3.
\]

\section{The Splitting Axiom}

The associativity of quantum multiplication follows from the splitting axiom, which describes how Gromov-Witten classes behave under gluing of stable curves.

Let $\varphi_S: \overline{\mathcal{M}}_{n_1+1} \times \overline{\mathcal{M}}_{n_2+1} \to \overline{\mathcal{M}}_n$ be the morphism associated with an ordered partition $S = (S_1, S_2)$ where $S_1 \coprod S_2 = \{1, \ldots, n\}$.

The splitting axiom states that for the Poincar\'e-dual class $\sum_{i,j} \eta^{i,j} \alpha_i \otimes \beta_j$ of the diagonal $\Delta \subset X \times X$:
\begin{equation}
\begin{aligned}
&\varphi_S^*(I_{n,d}^X(a_1 \otimes \cdots \otimes a_n)) \\
&= \sum_{d = d_1 + d_2} \sum_{i,j} I_{n_1+1, d_1}^X\left(\left(\bigotimes_{k_1 \in S_1} a_{k_1}\right) \otimes \alpha_i\right) \eta^{i,j} \otimes I_{n_2+1, d_2}^X\left(\beta_j \otimes \left(\bigotimes_{k_2 \in S_2} a_{k_2}\right)\right).
\end{aligned}
\tag{10}
\end{equation}

In particular, for $n=4$, the splitting axiom gives the associativity condition:
\begin{equation}
\begin{aligned}
&\sum_{d = d_1 + d_2} \sum_{i,j} I_{3, d_1}^X(a \otimes b \otimes \alpha_i) \eta^{i,j} I_{3, d_2}^X(\beta_j \otimes c \otimes d) \\
&= \sum_{d = d_1 + d_2} \sum_{i,j} I_{3, d_1}^X(b \otimes c \otimes \alpha_i) \eta^{i,j} I_{3, d_2}^X(\beta_j \otimes a \otimes d).
\end{aligned}
\tag{11}
\end{equation}
This identity is proved using the fact that $\overline{\mathcal{M}}_4 = \mathbb{CP}^1$.

\section{Equivariant Gromov-Witten Classes}

The equivariant version follows by replacing $X$ with $X_G = X \times_G EG$ and $BG$ with $BG$.

\begin{Definition}[\cite{Kim1996} Definition in Section 3.3]
\label{def:EquivariantGromovWittenClass}
Let $X$ have a continuous $G$-action with $G$ connected and compact. The equivariant Gromov-Witten class
\[
I_{n,d}^{X_G}: H_G^*(X)^{\otimes n} \to H^*(\overline{\mathcal{M}}_n) \otimes_{\mathbb{Q}} H^*(BG)
\]
is defined by
\begin{equation}
I_{n,d}^{X_G}(a_1 \otimes \cdots \otimes a_n) := \pi_*^{X_G} (ev^{X_G})^*(a_1 \otimes \cdots \otimes a_n),
\tag{12}
\end{equation}
where $a_i \in H_G^*(X)$.
\end{Definition}

The equivariant quantum cohomology module $H_G^*(X) \otimes_{\mathbb{Q}} \mathbb{Q}[q]$ has a unique multiplication characterized by:
\begin{equation}
\langle a_1 \cdot a_2, a_3 \rangle = \sum_d q^d I_{3,d}^{X_G}(a_1, a_2, a_3),
\tag{13}
\end{equation}
where $\langle, \rangle$ is the $q$-linear expansion of the equivariant Poincar\'e pairing.

The equivariant splitting axiom holds with all tensor products arising from $H^*(BG)$-module structures. The key observation is that all maps in the proof of the non-equivariant splitting property are (diagonally) equivariant.

\begin{Definition-Theorem}[\cite{Kim1996} Section 3.3]
\label{defth:EquivariantQuantumCohomology}
The equivariant quantum cohomology multiplication $QH_G(X)$ is defined analogously to the non-equivariant case. The associativity follows from the equivariant splitting property for $n=4$.
\end{Definition-Theorem}

When $G$ is trivial, the equivariant quantum cohomology reduces to ordinary quantum cohomology. Since $H^*(X)$ and $H^*(BG)$ are generated by even degree classes, $H_G^*(X) = H^*(X) \otimes_{\mathbb{Q}} H^*(BG)$ as linear spaces, and $QH_G^*(X)$ is a free $H^*(BG) \otimes_{\mathbb{Q}} \mathbb{Q}[q]$-module.

\section{Relation to Toda Lattice}

For complete flag manifolds $F_{1,2,\ldots,n}$, the polynomials $\Sigma_1(c, q), \ldots, \Sigma_n(c, q)$ turn out to be Poisson-commuting conservation laws of the Toda lattice. This remarkable connection, established in Givental-Kim \cite{GiventalKim1995}, shows that the quantum cohomology ring is isomorphic to the coordinate ring of the Toda lattice phase space.

The Toda lattice is a classical integrable system describing particles on a line with exponential interactions. Its phase space has a Poisson structure, and the quantities $\Sigma_i$ are independent integrals of motion that Poisson-commute. The quantum cohomology encodes this structure: the multiplication $\star$ on $H^*(F_{1,2,\ldots,n})$ corresponds to the Poisson bracket at the classical level, and the quantum parameters $q_i$ deform this structure.

The significance of this connection includes:
\begin{itemize}
\item It provides a computable system: the Toda lattice can be solved explicitly using algebraic methods
\item It connects enumerative geometry (counting rational curves) with classical mechanics
\item It suggests deep relationships between quantum cohomology and the Langlands program
\end{itemize}

Specifically, the classical limit ($q=0$) corresponds to the Poisson structure of the Gelfand-Dikii hierarchy, while the quantum deformation introduces the quantum parameters $q_i$ as deformation coordinates.

\section{Historical Remarks}

The results in this chapter represent a synthesis of several threads in algebraic geometry and symplectic topology:

\begin{itemize}
\item Gromov \cite{Gromov1985} introduced pseudo-holomorphic curves in symplectic manifolds
\item Floer \cite{Floer1986} developed quantum cohomology for symplectic manifolds
\item Kontsevich and Manin \cite{KontsevichManin1994} established the axioms of Gromov-Witten invariants
\item Astashkevich-Sadov \cite{AstashkevichSadov1995} and Kim \cite{Kim1995} independently conjectured the formula for flag manifolds
\item Ciocan-Fontanine \cite{CiocanFontanine1995} gave the first proof for complete flag manifolds $F_{1,2,\ldots,n}$ using a different method
\item Kim \cite{Kim1996} provided the systematic equivariant framework that works for partial flag manifolds
\item Lu also proved Theorem 2 independently using ``vertical quantum cohomology'' introduced by Astashkevich-Sadov \cite{Kim1996} Remark 2.
\end{itemize}

The connection to the Toda lattice, established in Givental-Kim \cite{GiventalKim1995}, is particularly significant: the quantum cohomology ring of the complete flag manifold is isomorphic to the coordinate ring of the Toda lattice phase space. This places quantum cohomology within the framework of integrable systems, where the polynomials $\Sigma_i(c,q)$ appear as Poisson-commuting conservation laws. This relationship has been further developed in the context of mirror symmetry and geometric Langlands correspondence.

The theory continues to develop through connections to integrable systems, mirror symmetry, and representation theory.

\endchapter*
